\exercice*

Commençons par résoudre séparément les deux inéquations.

\begin{minipage}[b]{.4\textwidth}
\begin{align*}
  (* for gauche, signe, droite, opération in etapes1 -*)
    (( gauche )) & (( signe )) (( droite )) \\
    (( gauche )) \textcolor{blue}{(( opération ))} & (* if loop.last *)(( solution1.symbole ))(* else *)(( signe ))(* endif *) (( droite )) \textcolor{blue}{(( opération ))}\\
  (* endfor *)
  x & (( solution1.symbole )) (( solution1.valeur|tex )) \\
\end{align*}
\end{minipage}%
\hfill%
\begin{minipage}[b]{.4\textwidth}
\begin{align*}
  (* for gauche, signe, droite, opération in etapes2 -*)
    (( gauche )) & (( signe )) (( droite )) \\
    (( gauche )) \textcolor{blue}{(( opération ))} & (* if loop.last *)(( solution2.symbole ))(* else *)(( signe ))(* endif *) (( droite )) \textcolor{blue}{(( opération ))}\\
  (* endfor *)
  x & (( solution2.symbole )) (( solution2.valeur|tex )) \\
\end{align*}
\end{minipage}%

Les solutions sont donc \textcolor{blue}{$x((solution1.symbole))((solution1.valeur|tex))$} (( etou )) \textcolor{Green}{$x((solution2.symbole))((solution2.valeur|tex))$}, dont voici la représentation sur la droite des réels.

\begin{center}
    \begin{tikzpicture}[line width=2pt, xscale=(* if (max-min)/5 >= 4 *)0.5(* elif (max-min)/5 == 3 *).75(* elif (max-min)/5 == 2 *)1(* else *)2(* endif *)]
      \draw[-Stealth] ( (( min-1 )), 0) -- ( (( max+1 )), 0);
      \foreach \x in { (( min )), (( min + 5 )), ..., (( max ))}{
          \draw (\x, -.2) -- (\x, .2);
          \path (\x, 0) node[above=5pt]{$\x$};
      }
      \foreach \x in { (( min )), ..., (( max ))}{
          \draw[line width=1.5pt] (\x, -.15) -- (\x, .15);
      }
      \begin{scope}[blue]
          \path ( (( solution1.valeur )), 0) node[below=9pt]{$(( solution1.valeur|tex ))$};
          \draw[arrows={Bracket[scale=1.5, inset=3pt,
                (*- if solution1.symbole in "<>" -*),reversed(*- endif -*)
               ]-},yshift=-2pt]
               ($
                  ( (( solution1.valeur )), 0)
                  +
                  (* if solution1.symbole in "<" *)
                      (4.5pt, 0)
                  (* elif solution1.symbole in ">" *)
                      (-4.5pt, 0)
                  (* elif solution1.symbole in "≤" *)
                      (.5pt, 0)
                  (* elif solution1.symbole in "≥" *)
                      (-.5pt, 0)
                  (* endif *)
               $)
               --
               (* if solution1.symbole in "<≤" *) ( (( min - 1 )), 0) (* else *) ( (( max + 0.6 )), 0) (* endif *)
               ;
      \end{scope}
      \begin{scope}[Green]
          \path ( (( solution2.valeur )), 0) node[below=9pt]{$(( solution2.valeur|tex ))$};
          \draw[arrows={Bracket[scale=1.5
                (*- if solution2.symbole in "<>" -*),reversed(*- endif -*)
               ]-},yshift=-4pt]
               ($
                  ( (( solution2.valeur )), 0)
                  +
                  (* if solution2.symbole in "<" *)
                      (4.5pt, 0)
                  (* elif solution2.symbole in ">" *)
                      (-4.5pt, 0)
                  (* elif solution2.symbole in "≤" *)
                      (.5pt, 0)
                  (* elif solution2.symbole in "≥" *)
                      (-.5pt, 0)
                  (* endif *)
               $)
               --
               (* if solution2.symbole in "<≤" *) ( (( min - 1 )), 0) (* else *) ( (( max + 0.6 )), 0) (* endif *)
               ;
      \end{scope}
  \end{tikzpicture}
\end{center}

Puisque nous voulons
(* if etou == "et" *) l'intersection (* else *) l'union (* endif *)
des deux solutions (car $x((solution1.symbole))((solution1.valeur|tex))$ \textbf{(( etou ))} $x((solution2.symbole))((solution2.valeur|tex))$), alors nous cherchons les parties de la droite des réels surlignées
(* if etou == "et" *) deux (* else *) au moins une (* endif *)
fois. Les solutions sont donc :
\[
    x \in ((solution|tex))
\]
